\begin{algorithme}
    \fonction{charToDirection}
        {c : Caractere}
        {Direction}
        {}
        {
            \cas{c}
            {
                {\casclausegenerale{'N'}{\retourner{NORD}}} 
                {\casclausegenerale{'E'}{\retourner{EST}}}  
                {\casclausegenerale{'S'}{\retourner{SUD}}} 
                {\casclausegenerale{'O'}{\retourner{OUEST}}}  
            }
        }
\end{algorithme}
\vspace{1cm}
\begin{algorithme}
    \fonction{determinerDirection}
        {caseDep : \naturel, caseAdj : \naturel, largeur : \naturel}
        {Direction}
        {}
        {
        \sialors{caseAdj = caseDep + 1} {\retourner{EST}}
        \sialors{caseAdj = caseDep - 1} {\retourner{OUEST}}
        \sialors{caseAdj = caseDep + largeur} {\retourner{SUD}}
        \sialors{caseAdj = caseDep - largeur} {\retourner{NORD}}
        
        \retourner{-1}
        }
\end{algorithme}
\vspace{1cm}
\begin{algorithme}
    \fonction{chargerVilleDepuisFichier}
        {nomFichier : Chaine}
        {\typePointeur{Ville}}
        {f : Fichier \\ ville : \typePointeur{Ville} \\ ligne : Chaine \\ etat, largeur, numDep, c1, c2 : \naturel \\ dirChar : Caractere \\ ptrC1, ptrC2 : \typePointeur{Case} \\ dir1, dir2 : Direction}
        {
            \affecter{f}{ouvrir(nomFichier, "r")}
            \sialors{f = \motclefPointeurNULL}
            {
                \ecrire{"Erreur lors de l'ouverture du fichier de ville"}
                \retourner{\motclefPointeurNULL}
            }

            \affecter{ville}{\motclefPointeurNULL}
            \affecter{etat}{0}

            \tantque{lireLigne(f, ligne)}
            {
                \sialors{estVide(ligne)} {\continuer}
                \sialors{ligne = "."} {\sortirBoucle}

                \sialorssinon{etat = 0}
                {
                    \sialorssinon{formatValide(ligne, "\%u", largeur) et largeur > 0}
                    {
                        \affecter{ville}{creerVille(largeur)}
                        \sialors{ville = \motclefPointeurNULL}
                        {
                            \instruction{fermer(f)}
                            \retourner{\motclefPointeurNULL}
                        }
                        \affecter{etat}{1}
                    }
                    {
                        \ecrire{"Erreur de format ou largeur invalide"}
                        \instruction{fermer(f)}
                        \retourner{\motclefPointeurNULL}
                    }
                }
                {
                    \sialorssinon{etat = 1}
                    {
                        \sialorssinon{formatValide(ligne, "\%u\%c", numDep, dirChar)}
                        {
                            \sialorssinon{obtenirCase(ville, numDep) $\neq$ \motclefPointeurNULL}
                            {
                                \instruction{fixerDepart(ville, numDep, charToDirection(dirChar))}
                                \affecter{etat}{2}
                            }
                            {
                                \ecrire{"Erreur : Case de départ invalide"}
                                \instruction{liberer(ville)}
                                \instruction{fermer(f)}
                                \retourner{\motclefPointeurNULL}
                            }
                        }
                        {
                            \ecrire{"Erreur de format pour le départ et la direction"}
                            \instruction{liberer(ville)}
                            \instruction{fermer(f)}
                            \retourner{\motclefPointeurNULL}
                        }
                    }
                    {
                        \sialorssinon{formatValide(ligne, "\%u-\%u", c1, c2)}
                        {
                            \affecter{ptrC1}{obtenirCase(ville, c1)}
                            \affecter{ptrC2}{obtenirCase(ville, c2)}
                            
                            \sialorssinon{ptrC1 $\neq$ \motclefPointeurNULL et ptrC2 $\neq$ \motclefPointeurNULL}
                            {
                                \affecter{dir1}{determinerDirection(c1, c2, largeur)}
                                \affecter{dir2}{determinerDirection(c2, c1, largeur)}

                                \sialorssinon{dir1 $\neq$ -1 et dir2 $\neq$ -1}
                                {
                                    \instruction{fixerCaseAdj(ptrC1, dir1, c2)}
                                    \instruction{fixerCaseAdj(ptrC2, dir2, c1)}
                                }
                                {
                                     \ecrire{"Avertissement : Passage non orthogonal ou non adjacent ignoré"}
                                }
                            }
                            {
                                \ecrire{"Avertissement: Passage entre cases non valides ignoré"}
                            }
                        } 
                        {
                            \sialorssinon{formatValide(ligne, "\%u", c1)}
                            {
                                \affecter{ptrC1}{obtenirCase(ville, c1)}
                                \sialorssinon{ptrC1 $\neq$ \motclefPointeurNULL}
                                {
                                    \instruction{fixerObligatoire(ptrC1, VRAI)}
                                }
                                {
                                    \ecrire{"Avertissement: Case obligatoire non valide ignorée"}
                                }
                            } 
                            {
                                \ecrire{"Avertissement: Ligne ignorée avec format non reconnu"}
                            }
                        }
                    }
                }
            }

            \instruction{fermer(f)}
            \retourner{ville}
        }


\begin{algorithme}
\procedure{genererOrdres}
{chemin : \typePointeur{Chemin}}
{nbMouv, i, d, diff : \naturel \\ dirCourante, dirCible : Direction \\ caseActu, caseSuiv : \typePointeur{Case} \\ routeGauche, routeDroite, estInter, dirTrouvee : \booleen}
{
    \affecter{nbMouv}{nbrCases(chemin)}
    \affecter{dirCourante}{obtenirDirDepart(chemin)}
    \ecrire{"AV"}

    \pour{i}{0}{nbMouv - 1}{}
    {
        \affecter{caseActu}{obtenirCaseAuChemin(chemin, i)}
        \affecter{caseSuiv}{obtenirCaseAuChemin(chemin, i + 1)}

        \sialors{caseActu \neq \motclefPointeurNULL \text{ \textbf{et} } caseSuiv \neq \motclefPointeurNULL}
        {
            \affecter{dirCible}{NORD}
            \affecter{dirTrouvee}{FAUX}

            \pour{d}{0}{3}{}
            {
                \sialors{non dirTrouvee \text{ \textbf{et} } obtenirVoisin(\valeurPointeur{caseActu}, d) = obtenirNumeroCase(caseSuiv)}
                {
                    \affecter{dirCible}{d}
                    \affecter{\text{dirTrouvee}}{VRAI}
                }
            }

            \sialors{dirTrouvee}
            {
                \affecter{routeGauche}{obtenirVoisin(\valeurPointeur{caseActu}, (dirCourante + 3) \% 4) \neq 0}
                \affecter{routeDroite}{obtenirVoisin(\valeurPointeur{caseActu}, (dirCourante + 1) \% 4) \neq 0}
                \affecter{estInter}{routeGauche \text{ \textbf{ou} } routeDroite}
                \affecter{diff}{(dirCible - dirCourante + 4) \% 4}

                \cas{diff}
                {
                    {\casclausegenerale{1}{\ecrire{"TD"} \ecrire{"AV"}}}
                    {\casclausegenerale{3}{\ecrire{"TG"} \ecrire{"AV"}}}
                    {\casclausegenerale{0}{\sialors{estInter}{\ecrire{"AV"}}}}
                }
                \affecter{dirCourante}{dirCible}
            }
        }
    }
    \ecrire{"."}
}
\end{algorithme}